% ===============================================================
% Section II  --  Dataset and Preprocessing
% Owner: Data Engineering Lead (integrated)
% ===============================================================
\section{Dataset and Preprocessing}
\label{sec:dataset}

We used the structured Shakespeare corpus provided for the assignment,
covering three plays: \emph{Hamlet}, \emph{Macbeth}, and \emph{Romeo and
Juliet}. The corpus is supplied in two granularities, a scene-level JSONL
file in which each line is one scene, and an utterance-level JSONL file in
which each line is a single speaker turn. We prepared both granularities
because they serve different retrieval needs: scene chunks preserve
narrative context, while utterances give precise, speaker-attributed
evidence. No external summaries or glossaries were introduced; all
enrichment metadata (scene summaries and keywords) was taken directly
from the provided dataset fields, so the report does not mix primary
Shakespearean text with outside explanatory material.

Table~\ref{tab:dataset-stats} summarises the corpus after preparation.

\begin{table}[ht]
\centering
\caption{Dataset statistics after preprocessing.}
\label{tab:dataset-stats}
\begin{tabular}{lccc}
\toprule
Play & Scenes & Scene chunks & Utterances \\
\midrule
Hamlet           & 20 & 63  & 1{,}147 \\
Macbeth          & 28 & 88  & 658   \\
Romeo and Juliet & 25 & 78  & 817   \\
\midrule
\textbf{Total}   & \textbf{73} & \textbf{229} & \textbf{2{,}622} \\
\bottomrule
\end{tabular}
\end{table}

The raw utterance files contained 3{,}613 records. Removing 991 stage
direction records (those whose \texttt{speaker} field was
\texttt{STAGE\_DIRECTION}) left 2{,}622 spoken utterances. Stage
directions describe physical action rather than dialogue and add noise to
retrieval, so they were excluded.

\subsection{Dataset Schema}
Each processed record follows the structure below, matching the format
recommended in the assignment specification:

\begin{lstlisting}
{
  "play":          "Macbeth",
  "act":           2,
  "scene":         2,
  "speaker":       null,
  "text":          "So foul and fair a day...",
  "source_id":     "macbeth_2_2",
  "scene_summary": "Macbeth murders Duncan; guilt begins.",
  "keywords":      ["murder", "guilt", "Duncan"]
}
\end{lstlisting}

For scene-level chunks the \texttt{speaker} field is \texttt{null}, since
a scene is kept whole rather than split by speaker. For utterance-level
chunks \texttt{speaker} holds the character name and \texttt{text} holds
only that character's line. The \texttt{source\_id} field lets any
retrieved chunk be traced back to its exact play, act, and scene.

\subsection{Cleaning and Filtering}
Three cleaning steps were applied to every record before indexing:
\begin{enumerate}
  \item \textbf{Stage-direction removal.} A regular expression stripped
  bracketed stage directions such as \texttt{[Exit Ghost.]} or
  \texttt{[Aside]}, and whole utterance records with speaker
  \texttt{STAGE\_DIRECTION} were dropped.
  \item \textbf{Whitespace normalisation.} Runs of spaces and newlines
  were collapsed to a single space, and leading/trailing whitespace was
  stripped from each text field.
  \item \textbf{Short-chunk removal.} Scene chunks under 50 characters
  and utterance chunks under 10 characters after cleaning were dropped,
  as such fragments carry too little meaning for reliable retrieval.
\end{enumerate}

\subsection{Chunking Strategy}
The retrieval unit was chosen separately for each index.

\textbf{Scene-level chunks} form the scene index. Each scene becomes one
chunk. A single line in isolation rarely makes sense to a beginner, the
line \emph{``To be or not to be''} says little without knowing who speaks
it and why, whereas a whole scene supplies enough context for a useful
answer. The trade-off is that scenes are long and may contain material
unrelated to a given question. To mitigate this, any scene longer than
400 words was split into overlapping chunks with a 50-word overlap so the
embedding model's 512-token limit is never silently exceeded. This
yielded \textbf{229 scene-level chunks}.

\textbf{Utterance-level chunks} form the utterance index. Each speaker
turn becomes its own chunk, giving precise retrieval and exact
attribution of who said what, in which scene. The trade-off is the loss
of surrounding context that scenes provide. The \textbf{2{,}622}
utterance chunks were stored in a ChromaDB vector collection, which also
supports filtering by play, act, scene, or speaker.

\textbf{Text enrichment.} For both granularities, the
\texttt{scene\_summary} and \texttt{keywords} were prepended to the chunk
text before embedding, so that retrieval operates on meaning rather than
surface word overlap. An enriched chunk looks like:

\begin{lstlisting}
[Scene: Guards and Horatio see the ghost
 and tell Hamlet.]
[Keywords: ghost, Horatio]
BERNARDO. Who's there? FRANCISCO. Nay,
 answer me...
\end{lstlisting}

This lets a question such as \emph{``Why does the ghost appear?''} match
the right scene even when the word \emph{appear} never occurs in the
dialogue. For utterance chunks the speaker name was also prepended so the
system can distinguish between characters. The chunking trade-off, short
units that are precise but context-poor versus long units that are
context-rich but less specific, is examined empirically in
Section~\ref{sec:results}.
